\documentclass[11pt,a4paper]{article}
\usepackage[margin=2.5cm]{geometry}
\usepackage[T1]{fontenc}
\usepackage{lmodern}
\usepackage{amsmath,amssymb,amsthm,mathtools,bm}
\usepackage{xcolor}
\usepackage[colorlinks=true,linkcolor=blue!50!black,citecolor=blue!50!black]{hyperref}

\newcommand{\m}{\bm{m}}
\newcommand{\h}{\bm{h}}
\newcommand{\vv}{\bm{v}}
\newcommand{\Hext}{\bm{H}}
\newcommand{\T}{\mathsf{T}}
\newcommand{\proj}{\mathsf{P}}
\newtheorem*{remark}{Remark}

\title{Learned coarse micromagnetics: the theory behind the torque and curvature method}
\author{Working notes}
\date{\today}

\begin{document}
\maketitle

\begin{abstract}
These notes collect the theory behind the learned coarse model: how a hysteresis loop is computed as a
sequence of constrained energy minimisations, how a learned energy term is trained from stored truth
states without ever relaxing during training, and why the second-order condition---the lowest curvature
of the energy on the product of unit spheres---is needed on top of the first-order (torque) condition.
They are notes for study, not a paper: the derivations are spelled out, results are only quoted where
they motivate a step.
\end{abstract}

\tableofcontents

%======================================================================
\section{Quasi-static hysteresis as a sequence of minimisations}
%======================================================================

A soft magnet driven slowly by an applied field $\Hext$ is, at every instant, in a local minimum of its
micromagnetic energy. The dynamics (precession, damping) only decides \emph{which} minimum it falls
into when the one it sits in disappears. A quasi-static hysteresis loop is therefore computed as a
sequence of minimisations: for field values $\Hext_0, \Hext_1, \dots$,
\begin{equation}
  \m_t = \operatorname{relax}\bigl(\m_{t-1};\, \Hext_t\bigr),
  \label{eq:loop}
\end{equation}
where $\operatorname{relax}$ descends the energy at the new field starting from the previous state. The
field history lives in the state: two different $\m_{t-1}$ at the same $\Hext_t$ can relax into
different minima. This is the whole content of hysteresis.

\subsection{Discrete energy and effective field}

On a grid of $N$ cubic cells of volume $V$ the magnetisation is $\m = (\m_1,\dots,\m_N)$ with
$|\m_i|=1$, so the state space is the product of spheres $\mathcal{M} = (S^2)^N$. The energy is
\begin{equation}
  E(\m;\Hext) = E_\mathrm{d}(\m) + E_\mathrm{ex}(\m) + E_\mathrm{Z}(\m;\Hext) + E_\theta(\m),
\end{equation}
and the effective field in cell $i$ is its negative gradient in field units,
\begin{equation}
  \h_i = -\frac{1}{\mu_0 M_s V}\,\frac{\partial E}{\partial \m_i}.
  \label{eq:heff}
\end{equation}
The three physical terms are
\begin{align}
  E_\mathrm{Z} &= -\mu_0 M_s V \sum_i \m_i\cdot\Hext,
    & \h^\mathrm{Z}_i &= \Hext, \\
  E_\mathrm{ex} &= \frac{A V}{\Delta x^2} \sum_{\langle ij\rangle} |\m_i-\m_j|^2,
    & \h^\mathrm{ex}_i &= \frac{2A}{\mu_0 M_s \Delta x^2} \sum_{j\sim i} (\m_j-\m_i), \\
  E_\mathrm{d} &= \frac{\mu_0 M_s^2 V}{2} \sum_{ij} \m_i\cdot \mathsf{N}_{ij}\m_j,
    & \h^\mathrm{d}_i &= -M_s \sum_j \mathsf{N}_{ij}\m_j,
\end{align}
with $\langle ij\rangle$ the nearest-neighbour pairs (a missing neighbour at the surface contributes
nothing: the free-surface condition), and $\mathsf{N}$ the demagnetisation tensor of cubic cells
(Newell), applied by FFT with zero padding, i.e.\ open boundaries. Exchange and demag are quadratic in
$\m$, so their fields are \emph{linear}: $\h^\mathrm{ex}+\h^\mathrm{d} = \mathsf{L}\m$ with
$\mathsf{L}$ symmetric and negative semidefinite. This will make their second derivatives trivial.

\subsection{The coarse grid and the learned energy}

The truth is simulated on a fine grid (10\,nm here, about two exchange lengths) and stored as block
averages. On the coarse grid (50--200\,nm) each coarse cell is the renormalised block average
\begin{equation}
  \m^{c}_I = \frac{\langle \m\rangle_I}{|\langle\m\rangle_I|},
  \qquad \langle\m\rangle_I = \frac{1}{|I|}\sum_{i\in I}\m_i .
\end{equation}
Exchange, demag and Zeeman are evaluated exactly on the coarse grid. What the coarse grid gets wrong
---mainly domain walls and vortex cores, which are much narrower than a coarse cell---is left to a
learned energy
\begin{equation}
  E_\theta(\m) = K_\mathrm{d} V \sum_I \rho_\theta(\m)_I, \qquad K_\mathrm{d} = \tfrac12\mu_0 M_s^2,
\end{equation}
where $\rho_\theta$ is a small CNN (three $3\times3\times3$ convolutions): a local energy density, so the
same weights apply to a grid of any size. Its field $\h^\theta = -\partial E_\theta/\partial\m /
(\mu_0 M_s V)$ comes from automatic differentiation. Learning an \emph{energy} rather than a field has
two consequences used below: the learned field is conservative, and its Jacobian is symmetric.

The last projection layer is initialised to zero, so $E_\theta\equiv0$ at the start and training begins
from plain coarse micromagnetics.

\begin{remark}[What renormalisation throws away]
$|\langle\m\rangle_I|<1$ wherever a wall or vortex passes through the block. At 50\,nm cells this is
negligible (median $0.999$, $0.2\,\%$ of cells below $0.9$); at 200\,nm it is not ($21\,\%$ of cells below
$0.9$). The length measures how much structure a cell hides; dropping it is one source of the closure
problem (section~\ref{sec:limits}).
\end{remark}

%======================================================================
\section{First order: equilibria and the torque loss}
%======================================================================

\subsection{Equilibrium on the product of spheres}

Minimising $E$ subject to $|\m_i|=1$ with Lagrange multipliers $\lambda_i$,
$\mathcal{L} = E + \sum_i \tfrac{\lambda_i}{2}(|\m_i|^2-1)$, gives
$\partial E/\partial\m_i + \lambda_i\m_i = 0$, i.e.\ $\h_i \parallel \m_i$. Equivalently the torque
vanishes, or the tangential part of the field does (Brown's condition):
\begin{equation}
  \m_i\times\h_i = 0
  \quad\Longleftrightarrow\quad
  \proj_i \h_i = 0, \qquad \proj_i = \mathsf{1}-\m_i\m_i^\T .
  \label{eq:brown}
\end{equation}
The multiplier is $\lambda_i = \mu_0 M_s V\,\m_i\cdot\h_i$: the field component along $\m_i$, which
reappears in the curvature (section~\ref{sec:curvature}).

\subsection{The torque loss}

Every stored truth state $\m^\star$ is a relaxed state of the fine physics at its field $\Hext^\star$.
If the coarse model were exact, its coarse average would be an equilibrium of the coarse energy. The
torque loss asks exactly that:
\begin{equation}
  L_\tau(\theta) = \frac{1}{N}\sum_i \frac{\bigl|\proj_i\,\h_i(\m^\star;\Hext^\star,\theta)\bigr|^2}{M_s^2},
  \label{eq:torque}
\end{equation}
averaged over the stored states. It is linear in $E_\theta$'s field, cheap (one forward and one
backward pass), and needs no relaxation: the data says where the equilibria are, and $E_\theta$ is
fitted so that they are equilibria.

\begin{remark}
An equilibrium can be a minimum, a saddle or a maximum. Equation~\eqref{eq:torque} cannot tell them
apart. A model trained on it alone can turn data states into saddles; the rollout~\eqref{eq:loop} then
rolls away from them at once. This is why the second-order condition is needed.
\end{remark}

%======================================================================
\section{Second order: the curvature on the sphere}
\label{sec:curvature}
%======================================================================

\subsection{\texorpdfstring{The Hessian on $(S^2)^N$}{The Hessian on the product of spheres}}

Let $\vv=(\vv_1,\dots,\vv_N)$ be a tangent vector at $\m$, $\m_i\cdot\vv_i=0$. The curve
\begin{equation}
  \m_i(t) = \frac{\m_i + t\vv_i}{|\m_i + t\vv_i|}
  = \m_i + t\vv_i - \tfrac{t^2}{2}|\vv_i|^2\m_i + O(t^3)
\end{equation}
stays on the spheres. Expanding the energy along it,
\begin{equation}
  \frac{d^2E}{dt^2}\Big|_{t=0}
  = \sum_{ij}\vv_i\cdot\frac{\partial^2E}{\partial\m_i\partial\m_j}\vv_j
    - \sum_i |\vv_i|^2\,\m_i\cdot\frac{\partial E}{\partial\m_i}.
\end{equation}
The second term is the curvature of the sphere itself: moving along a tangent direction drops the
point towards the centre, and the first-order gradient along $\m_i$ is felt at second order. With
$\partial E/\partial\m_i = -\mu_0M_sV\,\h_i$ and $\partial^2E/\partial\m_i\partial\m_j =
-\mu_0M_sV\,\partial\h_i/\partial\m_j$,
\begin{equation}
  \frac{1}{\mu_0M_sV}\frac{d^2E}{dt^2}\Big|_{t=0} = \vv\cdot\mathsf{A}\vv,
  \qquad
  (\mathsf{A}\vv)_i = (\m_i\cdot\h_i)\,\vv_i - \proj_i\bigl(D\h[\vv]\bigr)_i ,
  \label{eq:hessian}
\end{equation}
where $D\h[\vv] = \mathsf{L}\vv + D\h^\theta[\vv]$ is the directional derivative of the field (the
applied field does not depend on $\m$). $\mathsf{A}$, restricted to the tangent space, is the
Riemannian Hessian of $E$ on $(S^2)^N$ in units of A/m. It is symmetric because $\mathsf{L}$ is and
because $D\h^\theta = -\nabla^2E_\theta/(\mu_0M_sV)$ is a Hessian---one reason to learn an energy rather
than a field.

\subsection{\texorpdfstring{The lowest curvature $\kappa$}{The lowest curvature kappa}}

\begin{equation}
  \kappa(\m;\Hext) = \min_{\vv\in T_{\m},\ \vv\neq0} \frac{\vv\cdot\mathsf{A}\vv}{\vv\cdot\vv}
  \qquad [\mathrm{A/m}],
  \label{eq:kappa}
\end{equation}
the smallest eigenvalue of $\mathsf{A}$ on the tangent space: the curvature of the energy bowl along
its flattest direction. At an equilibrium,
\begin{itemize}
  \item $\kappa>0$: a strict local minimum, the state persists under the relaxation;
  \item $\kappa=0$: the bowl is flat in one direction---the branch of minima ends here;
  \item $\kappa<0$: a saddle; the minimiser rolls away along the softest direction.
\end{itemize}
We quote $\kappa/M_s$, which is dimensionless.

\paragraph{Reading the diagonal term.}
$(\m_i\cdot\h_i)$ is the field component along the magnetisation. For a cell aligned with a strong field
it is large and positive, a stiff spring; for a cell pointing against its local field it is negative
and drives instability. The applied field enters $\mathsf{A}$ only here, as $\m_i\cdot\Hext$. At a fixed
state, changing the field by $\Delta\Hext$ changes every diagonal entry by at most $|\Delta\Hext|$, so
$\kappa$ moves by at most $|\Delta\Hext|$ per field step (the states also move, which this bound does not
include). With $\mu_0|\Delta\Hext|=4$\,mT and $\mu_0M_s=1$\,T, one step moves $\kappa/M_s$ by about
$0.004$.

\subsection{Switching as a fold bifurcation}

Along a branch of minima $\m(\Hext)$, $\kappa>0$. The branch ends at a switching field $H_\mathrm{sw}$
where a minimum and a saddle merge (a saddle-node, or fold, bifurcation). Reduced to the amplitude $s$
of the softest mode, the energy near the fold has the normal form
\begin{equation}
  e(s;H) = -\alpha\,\varepsilon\, s + \tfrac{\beta}{3}s^3,\qquad \varepsilon = H_\mathrm{sw}-H,
  \quad \alpha,\beta>0 .
\end{equation}
For $\varepsilon>0$ the equilibria are $s_\pm=\pm\sqrt{\alpha\varepsilon/\beta}$, the minimum at $s_+$
with curvature $e''(s_+)=2\beta s_+ = 2\sqrt{\alpha\beta\,\varepsilon}$. Hence
\begin{equation}
  \kappa \propto \sqrt{H_\mathrm{sw}-H} \quad\text{as } H\to H_\mathrm{sw},
\end{equation}
and for $\varepsilon<0$ there is no equilibrium near $s=0$: the minimiser falls into another branch.
\textbf{A switching field is where $\kappa$ of the current branch reaches zero.} Getting switching
fields right therefore means getting $\kappa$ right near the end of each branch, not only getting the
states right; and $\kappa$ goes to zero steeply, as a square root, so a small error in the energy
landscape moves the switching field by a field step or more.

\subsection{The stability loss}

The torque loss makes data states equilibria; the stability loss makes them minima:
\begin{equation}
  L_\mathrm{stab}(\theta) = \Bigl\langle \operatorname{relu}\bigl(-\kappa(\m^\star;\Hext^\star,\theta)/M_s\bigr)^2 \Bigr\rangle,
  \qquad
  L = L_\tau + w\,L_\mathrm{stab}, \quad w=1 .
  \label{eq:stab}
\end{equation}
Only negative curvature is penalised. The states just before a switch are barely stable
($\kappa\gtrsim0$) and must stay so. A stronger weight pushes data states deeper into their minima and
distorts the energy ($w=10$ made the loop error 1.7 times larger); a margin $\kappa\ge\delta>0$ would
do the same to the pre-switch states and was not tried.

\begin{remark}[Upper and lower bounds on $\kappa$: what was tried]
A \emph{softness} condition $0\le\kappa/M_s\le0.004$ on the last state before each switch (stable, but
so little that one field step ends the branch) helped models trained for 3{,}000 steps, but not models
trained for 10{,}000; a \emph{switching} loss that pushed $\kappa$ of pinned states below $-0.005$
contradicted the stability loss, because the pinned state is next to a data state that must have
$\kappa\ge0$, and one field step moves $\kappa$ only by $\sim0.004$. Only the lower bound~\eqref{eq:stab}
survived.
\end{remark}

%======================================================================
\section{\texorpdfstring{Computing $\kappa$ and its gradient}{Computing kappa and its gradient}}
%======================================================================

\subsection{Hessian-vector products}

$\mathsf{A}$ is never formed: it has $(3N)^2$ entries. Only products $\mathsf{A}\vv$ are needed, and
each costs about one field evaluation:
\begin{itemize}
  \item $\mathsf{L}\vv$: demag and exchange are linear, so $\mathsf{L}\vv = \h^\mathrm{d}(\vv)+\h^\mathrm{ex}(\vv)$,
        the same FFT convolution and stencil applied to $\vv$;
  \item $D\h^\theta[\vv]$: one vector-Jacobian product of $\h^\theta(\m)$ with $\vv$ by automatic
        differentiation; because the Jacobian is symmetric, the vector-Jacobian product equals the
        Jacobian-vector product. In code:
        \texttt{torch.autograd.grad(h\_theta, m, grad\_outputs=v)};
  \item $(\m_i\cdot\h_i)\vv_i$ and the projections $\proj_i$: pointwise.
\end{itemize}
Since $\h^\theta$ is itself a gradient of the CNN, $D\h^\theta[\vv]$ is a second derivative of the CNN,
and the gradient of the loss with respect to $\theta$ a third. The activation therefore has to be smooth
(SiLU or GELU, not ReLU, whose second derivative vanishes almost everywhere).

\subsection{Lanczos for the lowest eigenvalue}

With these products, the Lanczos method builds an orthonormal basis $Q_k=[\bm q_1,\dots,\bm q_k]$ of the
Krylov space $\operatorname{span}\{\vv_0,\mathsf{A}\vv_0,\dots,\mathsf{A}^{k-1}\vv_0\}$ from a random
tangent start $\vv_0$, and the small tridiagonal matrix $T_k = Q_k^\T\mathsf{A}Q_k$. The lowest
eigenvalue of $T_k$ (a Ritz value) converges quickly to the lowest eigenvalue of $\mathsf{A}$, because
it lies at the edge of the spectrum; its eigenvector $\bm y$ gives the Ritz vector $\vv = Q_k\bm y$. We use
$k=40$ with full reorthogonalisation (twice, for float32), which gives $\kappa$ to a few per cent when it
is negative, the case the loss cares about.

The Ritz value is a Rayleigh quotient, so it is an \emph{upper} bound on the true $\kappa$. The
stability loss can therefore miss an instability that Lanczos has not resolved, but never penalises a
true minimum.

\subsection{\texorpdfstring{The gradient of $\kappa$ with respect to $\theta$}{The gradient of kappa with respect to theta}}

Training needs $\partial\kappa/\partial\theta$. For a simple eigenvalue with normalised eigenvector
$\vv$, the derivative of $\kappa = \vv\cdot\mathsf{A}_\theta\vv$ is
\begin{equation}
  \frac{\partial\kappa}{\partial\theta}
  = \vv\cdot\frac{\partial\mathsf{A}_\theta}{\partial\theta}\vv
  + 2\,\frac{\partial\vv}{\partial\theta}\cdot\mathsf{A}_\theta\vv
  = \vv\cdot\frac{\partial\mathsf{A}_\theta}{\partial\theta}\vv ,
\end{equation}
since $\mathsf{A}_\theta\vv=\kappa\vv$ and $\vv\cdot\partial\vv/\partial\theta=0$ (Hellmann--Feynman).
So the eigenvector can be held fixed: Lanczos runs without gradients, and one more Hessian-vector product
$\vv\cdot\mathsf{A}_\theta\vv/\vv\cdot\vv$, with gradients on, gives a differentiable $\kappa$. The data
state $\m^\star$ is fixed as well; only $\theta$ moves.

%======================================================================
\section{Relaxation: the rollout}
%======================================================================

At test time a loop is computed by~\eqref{eq:loop}. The relaxation descends the energy on the spheres
with a norm-preserving update (magnum.np's midpoint rule, a Cayley-type step in the plane of $\m_i$
and $\m_i\times(\m_i\times\h_i)$) and Barzilai--Borwein step sizes: with $\bm s = \m^{(k+1)}-\m^{(k)}$ and
$\bm y$ the change in the torque $\m\times(\m\times\h)$, the step alternates between
$\tau = \bm s\cdot\bm s/\bm s\cdot\bm y$ and $\bm s\cdot\bm y/\bm y\cdot\bm y$. It stops when the largest
torque in any cell is below 10\,A/m. The relaxation is not part of training; it is how the model is
used and evaluated.

The loop error is the mean over field steps of $|\langle\m\rangle_\mathrm{model}-\langle\m\rangle_\mathrm{truth}|$
along the field axis, starting both from the first truth state.

%======================================================================
\section{What the losses do not constrain}
\label{sec:limits}
%======================================================================

\begin{itemize}
  \item \textbf{The landscape away from the data.} Torque and stability are conditions at the stored
        states only. Branch ends (where $\kappa\to0$) and minima nowhere near the data are left free. A
        model trained long enough can build a deep spurious minimum: one run fell into a demagnetised state
        that no field up to 200\,mT could leave, with training losses that looked normal. Choosing the
        checkpoint by its rollout of the validation loops guards against this.
  \item \textbf{Closure.} Many fine states share one coarse $\m$: the position of a wall inside a cell,
        its chirality, a vortex's polarity. They have different energies and different futures, and an
        energy of the coarse $\m$ alone has to average over them. Hidden variables per cell (an energy
        $E_\theta(\m,\bm z)$ minimised over both) or the block-average length are ways to carry that
        information.
  \item \textbf{Unseen cells.} The energy density is local and the same everywhere. A model trained on
        plates two cells thick has only seen surface cells, and fails in the interior of a thick block.
\end{itemize}

%======================================================================
\begin{thebibliography}{9}
\bibitem{brown} W.\,F.~Brown, \emph{Micromagnetics} (Interscience, 1963).
\bibitem{absil} P.-A.~Absil, R.~Mahony, R.~Sepulchre, \emph{Optimization Algorithms on Matrix
  Manifolds} (Princeton University Press, 2008): Riemannian gradient and Hessian on spheres.
\bibitem{bb} J.~Barzilai, J.\,M.~Borwein, Two-point step size gradient methods, IMA J.\ Numer.\ Anal.\
  \textbf{8}, 141 (1988).
\bibitem{exl} L.~Exl et al., LaBonte's method revisited: an effective steepest descent method for
  micromagnetic energy minimization, J.\ Appl.\ Phys.\ \textbf{115}, 17D118 (2014).
\bibitem{lanczos} C.~Lanczos, An iteration method for the solution of the eigenvalue problem of
  linear differential and integral operators, J.\ Res.\ Natl.\ Bur.\ Stand.\ \textbf{45}, 255 (1950).
\bibitem{magnumnp} F.~Bruckner et al., magnum.np: a PyTorch based GPU enhanced finite difference
  micromagnetic simulation framework for high level development and inverse design, Sci.\ Rep.\
  \textbf{13}, 12054 (2023).
\end{thebibliography}

\end{document}
